\setcounter{chapter}{11} % Sets chapter counter so the current chapter is 12
\renewcommand{\thetheorem}{\thechapter.\arabic{theorem}}
\setcounter{theorem}{0}

\chapter{Dimension and Subspaces}
In this chapter, we synthesize the fundamental results regarding finite-di\-men\-sional vector spaces and explore the relationship between a space and its subspaces through the lens of dimensionality.

\begin{summary}[Properties of Finite-Dimensional Spaces]
Let $V$ be a finite-di\-men\-sional vector space over a field $K$. The following structural principles hold:
\begin{enumerate}
    \item Every finite list of vectors that spans $V$ contains a sublist which constitutes a \qt{basis} of $V$. Formally, if $(v_1, \dots, v_n)$ spans $V$, then there exists a sublist $(v_{i_1}, \dots, v_{i_d})$ with $1 \le i_1 < i_2 < \dots < i_d \le n$ that is a basis for $V$, where $d = \dim V$.
    \item Every linearly independent list of vectors in $V$ can be extended to a basis of $V$. If the list $(u_1, \dots, u_m)$ is linearly independent, there exist vectors $w_1, \dots, w_{d-m}$ such that 
    \[(u_1, \dots, u_m, w_1, \dots, w_{d-m})\] 
    is a basis of $V$, where $d = \dim V$.
\end{enumerate}
\end{summary}

% thm:12.1
\begin{theorem}[Equivalent Characterizations of a Basis]
\label{thm:basis_equivalent_characterizations}
Let $V$ be a vector space with $\dim V = n \in \mathbb{Z}_{\ge 0}$. Let $v_1, \dots, v_n \in V$ be a list of $n$ vectors. Then the following statements are equivalent:
\begin{enumerate}
    \item[\textbf{(a)}] The vectors $v_1, \dots, v_n$ are linearly independent.
    \item[\textbf{(b)}] The vectors $v_1, \dots, v_n$ generate $V$, i.e., $\Sp(v_1, \dots, v_n) = V$.
    \item[\textbf{(c)}] The vectors $v_1, \dots, v_n$ constitute a basis for $V$.
\end{enumerate}
\end{theorem}

\begin{proof}
The equivalence is established through the following implications:

\textbf{(a) $\implies$ (c):} Suppose that \textbf{(a)} holds, so the list $v_1, \dots, v_n$ is linearly independent. According to the property that every linearly independent list can be extended to a basis, and observing that the list already possesses length $n = \dim V$, the extension requires no additional vectors. Thus, the list is already a basis for $V$, satisfying \textbf{(c)}.

\textbf{(c) $\implies$ (a) \& (b):} This follows directly from the definition of a basis, which by requirement must satisfy both the condition of linear independence in \textbf{(a)} and the spanning property in \textbf{(b)}.

\textbf{(b) $\implies$ (c):} If the spanning property in \textbf{(b)} holds, the list $v_1, \dots, v_n$ generates $V$. By the properties of finite-dimensional spaces, this list must contain a sublist that is a basis. Since $\dim V = n$, any basis must have exactly length $n$. Therefore, the sublist must be the entire list itself, making $v_1, \dots, v_n$ a basis as required by \textbf{(c)}.
\end{proof}



% prop:12.2
\begin{proposition}[Dimensional Constraints on Subspaces]
\label{prop:subspace_dimension_constraints}
Let $V$ be a finite-di\-men\-sional vector space over $K$. Then every linear subspace $U \subseteq V$ is also finite-di\-men\-sional. Moreover, $\dim U \le \dim V$, with the equality $\dim U = \dim V$ holding if and only if $U = V$.
\end{proposition}

\begin{proof}
Let $n = \dim V$. We first establish that $U$ possesses a finite basis. If $U = \{0\}$, the empty set serves as the basis and we are done. Assume $U \neq \{0\}$. We select a non-zero vector $v_1 \in U$. If $\Sp(v_1) = U$, we have found our basis. 

If $\Sp(v_1) \neq U$, we choose a second vector $v_2 \in U \setminus \Sp(v_1)$. It follows from the properties of linear independence that the list $\{v_1, v_2\}$ is linearly independent. We proceed iteratively: after $j$ steps, we obtain a linearly independent list $\{v_1, \dots, v_j\} \subseteq U$. If $\Sp(v_1, \dots, v_j) = U$, the process terminates. Otherwise, we choose $v_{j+1} \in U \setminus \Sp(v_1, \dots, v_j)$. By the principle that adding a vector outside the span of a linearly independent set preserves independence, the expanded list $\{v_1, \dots, v_{j+1}\}$ remains linearly independent.

Crucially, in a space $V$ of dimension $n$, no linearly independent list can exceed length $n$. Consequently, this algorithmic procedure must terminate in $k$ steps, where $k \le n$. At this stage, we have a list $\{v_1, \dots, v_k\}$ that is linearly independent and spans $U$, thereby constituting a basis for $U$. This implies $U$ is finite-dimensional and $\dim U = k \le n = \dim V$.

Furthermore, if $\dim U = \dim V = n$, then $U$ contains a linearly independent list of length $n$. By Theorem \ref{thm:basis_equivalent_characterizations}, such a list satisfies property \textbf{(a)} and is therefore also a basis for $V$. This implies $\Sp(v_1, \dots, v_n) = V$. Since $U = \Sp(v_1, \dots, v_n)$ by the construction of the basis, we conclude $U = V$.
\end{proof}